\section{Глосарій}
\label{sec:hlosarii}

У \tabref{tab:hlosarii} наведено визначення ключових термінів, використаних у цьому документі.

\begin{table}[h]
\caption{Глосарій системи}
\label{tab:hlosarii}
\begin{tabular}{L{3.5cm} L{3.5cm} L{7cm}}
\toprule
\textbf{Термін (укр.)} & \textbf{Англ. відповідник} & \textbf{Визначення} \\
\midrule
Каталог & Catalog & Повний список усіх елементів контенту, доданих користувачем. \sked{K1} \\
\addlinespace
Елемент контенту & Content Item & Одна позиція в каталозі~--- книга, серіал або фільм із власними атрибутами (тип, статус, оцінка, теги тощо). \sked{K2} \\
\addlinespace
Добірка & Collection & Створена користувачем іменована група довільних елементів каталогу; один елемент може входити до кількох добірок одночасно. \sked{K3} \\
\addlinespace
Session & Session & Окремий запис факту перегляду або прочитання елемента; містить дату, оцінку та нотатку; для одного елемента може бути кілька Sessions. \sked{K4} \\
\addlinespace
Тег & Tag & Довільна текстова мітка, яку користувач додає до елемента для категоризації. \sked{K5} \\
\addlinespace
Статус & Status & Поточний стан роботи з елементом: \emph{заплановано~/ у процесі~/ завершено~/ відкладено}. \sked{K6} \\
\bottomrule
\end{tabular}
\end{table}
